\SourceLeaf{270}
\Article{17}
Resumamus ad maiorem illustrationem exemplum art. 13, ubi invenimus
\[
\begin{aligned}
A'=A'''&=2+4R+R^3+2R^4=2R-2R^2-R^3\\
A''&=2+R+4R^2+2R^3=-R+2R^2-2R^4
\end{aligned}
\]
Accipiendo pro $R$ valorem $\cos72^0+i\sin72^0$, erit
\[
\begin{aligned}
A'=A'''&=2\cos72^0-3\cos144^0+i(2\sin72^0-\sin144^0)\\
A''&=-3\cos72^0+2\cos144^0+i(\sin72^0+2\sin144^0)
\end{aligned}
\]
Determinabuntur itaque anguli $b',b''$ per aequationes
\[
\begin{aligned}
1)\qquad\sin b'&=\frac{2\sin72^0-\sin144^0}{\surd11}\\
2)\qquad\cos b'&=\frac{2\cos72^0-3\cos144^0}{\surd11}\\
3)\qquad\operatorname{tang}b'&=\frac{2\sin72^0-\sin144^0}{2\cos72^0-3\cos144^0}\\
4)\qquad\sin b''&=\frac{\sin72^0+2\sin144^0}{\surd11}\\
5)\qquad\cos b''&=\frac{-3\cos72^0+2\cos144^0}{\surd11}\\
6)\qquad\operatorname{tang}b''&=\frac{\sin72^0+2\sin144^0}{-3\cos72^0+2\cos144^0}
\end{aligned}
\]
Quaelibet aequationum $1,2,3$ sufficit ad determinandum angulum $b'$, si quadrans in quo accipiendus est innotuerit; hoc e signis quantitatum $2\sin72^0-\sin144^0$, $2\cos72^0-3\cos144^0$ decidi debebit: idem valet de angulo $b''$. In casu nostro $b'$ accipietur inter $0$ et $90^0$, $b''$ inter $90^0$ et $180^0$. Si aequationis $3$ numerator et denominator multiplicantur per $-3\cos72^0+2\cos144^0$, transibit in hanc
\SourceLeaf{271}
\[
\operatorname{tang}b'=\frac2{31}\{-\sin72^0+13\sin144^0\}
\]
et perinde ex aequatione $6$, multiplicato numeratore et denominatore per $2\cos72^0-3\cos144^0$, prodit
\[
\operatorname{tang}b''=\frac2{31}\{-13\sin72^0-\sin144^0\}
\]
Hinc fit in numeris
\[
\begin{aligned}
\operatorname{tang}b'&=+0{,}4316226944,\quad\log\operatorname{tang}b'=9{,}6351042715\\
&\hspace{1em}b'=23^020'46''04603\\
\operatorname{tang}b''&=-0{,}8355819332,\quad\log\operatorname{tang}b''=9{,}9219890411\mathrm n\\
&\hspace{1em}b''=140^07'6''52441
\end{aligned}
\]
unde derivatur
\[
5f'=186^048'38''61647,\qquad f'=37^021'43''723294
\]

.Habemus itaque
\[
\begin{aligned}
(2,1)&=-\frac15+\frac{\surd11}5\{2\cos37^021'43''723294+2\cos51^022'41''400558\}\\
(2,2)&=-\frac15+\frac{\surd11}5\{2\cos325^021'43''723294+2\cos267^022'41''400558\}\\
(2,4)&=-\frac15+\frac{\surd11}5\{2\cos253^021'43''723294+2\cos123^022'41''400558\}\\
(2,8)&=-\frac15+\frac{\surd11}5\{2\cos181^021'43''723294+2\cos339^022'41''400558\}\\
(2,5)&=-\frac15+\frac{\surd11}5\{2\cos109^021'43''723294+2\cos195^022'41''400558\}
\end{aligned}
\]
unde invenitur
\[
\begin{array}{rclcl}
(2,1)&=&+1{,}6825070652&=&2\cos\frac{360^0}{11}\\[3pt]
(2,2)&=&+0{,}8308299&=&2\cos\frac{720^0}{11}\\[3pt]
(2,4)&=&&=&2\cos\frac{1440^0}{11}\\[3pt]
(2,8)&=&&=&2\cos\frac{2880^0}{11}\\[3pt]
(2,5)&=&&=&2\cos\frac{1800^0}{11}
\end{array}
\]

\Article{18}
Exemplum aliud nobis suppeditabit aequatio $x^{17}-1=0$, quam per aliam methodum iam in \emph{Disquis. Arithm.} pertractaveramus. Statuemus itaque $n=17$, $\beta=8$, $g=3$; hinc respondent
\SourceLeaf{272}
\[
\begin{array}{r*{16}{@{\;}r}}
\text{numeris}&1\pdot&2\pdot&3\pdot&4\pdot&5\pdot&6\pdot&7\pdot&8\pdot&9\pdot&10\pdot&11\pdot&12\pdot&13\pdot&14\pdot&15\pdot&16\\
\text{indices}&0\pdot&14\pdot&1\pdot&12\pdot&5\pdot&15\pdot&11\pdot&10\pdot&2\pdot&3\pdot&7\pdot&13\pdot&4\pdot&9\pdot&6\pdot&8
\end{array}
\]
Hinc invenimus
\[
\begin{aligned}
A'=A''''''&=2R+2R^2+3R^4+4R^5+2R^6+2R^7\\
A''=A'''''&=2+3R+R^3+R^4+3R^5+4R^6+R^7\\
A'''=A''''&=3+3R+2R^2+3R^3+R^5+2R^6+R^7
\end{aligned}
\]
sive, quum in hoc casu fiat $R^4+1=0$
\[
\begin{aligned}
A'=A''''''&=-3-2R-2R^3\\
A''=A'''''&=1-4R^2\\
A'''=A''''&=3+2R+2R^3
\end{aligned}
\]
Statuendo itaque $R=\cos45^0+i\sin45^0$ erit
\[
\begin{aligned}
A'=A''''''&=-3-2i\surd2,\\
A''=A'''''&=1-4i,\\
A'''=A''''&=3+2i\surd2
\end{aligned}
\]

Invenientur itaque $b',b'',b'''$ per aequationes
\[
\begin{aligned}
\sin b'&=-\surd\tfrac8{17}&\sin b''&=-\surd\tfrac{16}{17}&\sin b'''&=+\surd\tfrac8{17}\\
\cos b'&=-\surd\tfrac9{17}&\cos b''&=+\surd\tfrac1{17}&\cos b'''&=+\surd\tfrac9{17}\\
\operatorname{tang}b'&=+\surd\tfrac89&\operatorname{tang}b''&=-4&\operatorname{tang}b'''&=+\surd\tfrac89
\end{aligned}
\]
unde deducimus
\[
b'=223^018'49'',\quad b''=284^02'10,\quad b'''=43^018'49''=b'-180^0
\]
\[
4f'=550^039'48'',\qquad f'=137^039'57''
\]
\[
\begin{aligned}
(2,1)&=-\frac18+\frac{\surd17}8\{2\cos137^039'57''+2\cos52^01'5''+2\cos265^038'52''+1\}\\
(2,4)&=-\frac18+\frac{\surd17}8\{2\cos92^039'57''+2\cos322^01'5''+2\cos130^038'52''-1\}\\
(2,9)&=-\frac18+\frac{\surd17}8\{2\cos47^039'57''+2\cos232^01'5''+2\cos355^038'52''+1\}\\
(2,10)&=-\frac18+\frac{\surd17}8\{2\cos2^039'57''+2\cos142^01'5''+2\cos220^038'52''-1\}\\
(2,13)&=-\frac18+\frac{\surd17}8\{2\cos317^039'57''+2\cos52^01'5''+2\cos85^038'52''+1\}\\
(2,5)&=-\frac18+\frac{\surd17}8\{2\cos272^039'57''+2\cos322^01'5''+2\cos310^038'52''-1\}\\
(2,15)&=-\frac18+\frac{\surd17}8\{2\cos227^039'57''+2\cos232^01'5''+2\cos175^038'52''+1\}\\
(2,11)&=-\frac18+\frac{\surd17}8\{2\cos182^039'57''+2\cos142^01'5''+2\cos40^038'52''-1\}
\end{aligned}
\]

\SourceLeaf{273}
\[
\begin{array}{rclcl}
\frac12(2,1)&=&+0{,}092268&=&\cos\frac4{17}360^0\\[3pt]
\frac12(2,3)&=&&=&\cos\frac5{17}360^0\\[3pt]
\frac12(2,9)&=&&=&\cos\frac2{17}360^0\\[3pt]
\frac12(2,10)&=&&=&\cos\frac6{17}360^0\\[3pt]
\frac12(2,13)&=&+0{,}93247&=&\cos\frac1{17}360^0\\[3pt]
\frac12(2,5)&=&&=&\cos\frac3{17}360^0\\[3pt]
\frac12(2,15)&=&&=&\cos\frac8{17}360^0\\[3pt]
\frac12(2,11)&=&&=&\cos\frac7{17}360^0
\end{array}
\]
\begin{center}*\qquad\qquad*\\*\end{center}

Ab his disquisitionibus generalioribus supra functiones $[r,R]$, quae theoriam secundam aequationum purarum in art. 360 \emph{Disquiss. Ar.} inchoatam magis illustrant et ampliant, ad casuum quorundam specialium considerationem accuratiorem (puta si pro $\beta$ valores determinati accipiuntur) progredimur; plures hinc investigationes non minus fertiles quam elegantes prodibunt, quarum aliae quidem iam in \emph{Disq. Ar.} (artt.\ .\ .\ .) pertractatae erant (sed per methodum diversam), aliae vero tamquam prorsus novae considerandae sunt. Mirum vero nexum inter hasce disquisitiones Arithmeticamque sublimiorem, quae incrementa maxima hactenusque inexspectata inde capit, in commentatione alia mox publici iuris facienda evolvere nobis reservamus. — Ceterum in tota disquisitione sequente supponemus, pro $r$ accipi radicem \emph{propriam} aequationis $x^n-1=0$, et pro $R$ radicem propriam aequationis $R^\beta-1=0$.

\Article{19}
Initium facimus a valore $\beta=2$, ubi itaque pro $R$ accipiendus est valor $-1$. Functio itaque nostra $[r,R]$ fit
\[
=r-r^g+r^{g^2}-r^{g^3}\dotsfour-r^{g^{n-2}}
\]
habeturque
\[
[r,R]=-[r^g,R]=+[r^{g^2},R]=-[r^{g^3},R]\quad\text{etc.}
\]
et generaliter, designante $\lambda$ integrum quemcunque per $n$ non divisibilem
\SourceLeaf{274}
\[
\begin{aligned}
[r^\lambda,R]&=+[r,R]\quad\text{si }\lambda\text{ est residuum quadraticum ipsius }n,\\
\relax[r^\lambda,R]&=-[r,R]\quad\text{si }\lambda\text{ est non-residuum quadraticum ipsius }n.
\end{aligned}
\]
Porro patet, si residua quadratica ipsius $n$ inter $1,2,3\dots n-1$ contenta indefinite designentur per $a$, atque non-residua ipsius $n$ inter eosdem limites per $b$, numeros
\[
1,g^2,g^4\dotsfour g^{n-3}
\]
si ad ordinem non respiciatur, congruos esse secundum modulum $n$ numeris $a$, et perinde numeros
\[
g,g^3,g^5\dotsfour g^{n-2}
\]
congruos ipsis $b$, ita ut fiat $[r,R]=\Sigma r^a-\Sigma r^b$.

Quodsi itaque statuimus $\frac{360^0}n=\omega$, atque $r=\cos k\omega+i\sin k\omega$, erit $[r,R]=\Sigma\cos ak\omega-\Sigma\cos bk\omega+i\Sigma\sin ak\omega-i\Sigma\sin bk\omega$. Iam per art. 14 quadratum functionis $[r,R]$ erit $=+n$ vel $=-n$, prout $n$ est formae $4z+1$ vel $4z-1$, adeoque in casu priori $[r,R]=\pm\surd n$, in posteriori $[r.R]=\pm i\surd n$; signum vero quantitati radicali praefixum ambiguum manet. Hinc derivantur summationes sequentes

\noindent I.\quad Si $n$ est formae $4z+1$
\[
\begin{aligned}
\Sigma\cos ak\omega-\Sigma\cos bk\omega&=\pm\surd n\\
\Sigma\sin ak\omega-\Sigma\sin bk\omega&=0
\end{aligned}
\]
\noindent II.\quad Si $n$ est formae $4z-1$
\[
\begin{aligned}
\Sigma\cos ak\omega-\Sigma\cos bk\omega&=0\\
\Sigma\sin ak\omega-\Sigma\sin bk\omega&=\pm\surd n
\end{aligned}
\]

Praeterea quum manifesto totus complexus numerorum $a,b$ conveniat cum his $1,2,3\dots n-1$, fit $\Sigma r^a+\Sigma r^b=r+r^2+r^3+\text{etc.}+r^{n-1}=-1$, et proin $\Sigma\cos ak\omega+\Sigma\cos bk\omega=-1$, $\Sigma\sin ak\omega+\Sigma\sin bk\omega=0$. Hinc e summationibus praecedentibus demanant sequentes:

\noindent I.\quad Pro casu priori
\[
\begin{aligned}
\Sigma\cos ak\omega&=-\tfrac12\pm\tfrac12\surd n\\
\Sigma\cos bk\omega&=-\tfrac12\mp\tfrac12\surd n\\
\Sigma\sin ak\omega&=\Sigma\sin bk\omega=0
\end{aligned}
\]

\SourceLeaf{275}
\noindent II.\quad Pro casu posteriori
\[
\begin{aligned}
\Sigma\cos ak\omega&=\Sigma\cos bk\omega=-\tfrac12\\
\Sigma\sin ak\omega&=\pm\tfrac12\surd n\\
\Sigma\sin bk\omega&=\mp\tfrac12\surd n
\end{aligned}
\]
Hae summationes per methodum haud multum diversam in \emph{Disquiss. Arr.} art. 356 iam sunt erutae; neutra quidem methodus ambiguitatem signi quantitati radicali praefigendi tollere valet, attamen hunc defectum in commentatione peculiari nuper supplevimus, ubi demonstratum est, pro valore $k=1$ signa superiora in omnibus formulis allatis accipi debere.
\BookRule{.3\textwidth}
